% page 470 du fac-simile (image p-469 du scan)
% bloc 157.5 x 242.0 mm ; marge gauche 8.0 mm
\pgc{-12.700mm}{0.000mm}{
\l{\cel{3}462}
\l{}
\l{\sou{procesionar}. (\sou{netrans}.)I(\sou{aludante la klerikaro od amaso}\cel{1}de}
\l{\cel{3}\sou{personi qui defilas ordinoze}). Marchar kolone en la kirko}
\l{\cel{3}od extere, kantante psalmi, himni, litanii, okazione di}
\l{\cel{3}ula festi religiala. - II. (\sou{metaf}.)(\sou{aludante longa serio}}
\l{\cel{3}de\cel{1}\sou{personi qui avancas ganso-marche}). Marchar ica dop ita,}
\l{\cel{3}unope. - DEFIRS.}
\l{}
\l{\sou{procidar}. (\sou{netrans.)}\cel{2}(\sou{teol.)}\cel{2}Exekutar la ago per qua la San-}
\l{\cel{3}ta Spirito pasas de la Patro a la Filio (en la religio ka-}
\l{\cel{3}tolika). -}
\l{}
\l{\sou{proda}. Qua agis plu kam brave. FIS.}
\l{}
\l{prodigar. (\sou{trans.)}\cel{2}I. Spensar sen modereso, sen limito.-}
\l{\cel{3}II. (\sou{metaf.)}\cel{2}Donacar sen modereso, sen limito. - EF.}
\l{}
\l{prodromo.\cel{2}(\sou{medicino)}.\cel{2}Signo qua anuncas morbo, maladesko.}
\l{\cel{3}- DEFIS.}
\l{}
\l{\sou{produktar}. (\sou{trans.)}\cel{2}Igar existar, igar naskar. - DEFIRS.}
\l{}
\l{\sou{produto}. (\sou{matem.}) Rezultajo de la multipliko di nombro per}
\l{\cel{3}altra. - F.}
\l{}
\l{\sou{profana}. I. Qua esas exter la kozi sakra. - II. (\sou{metaf.})}
\l{\cel{3}Qua esas exter ula cienco, ula arto. - DEFIS.}
\l{}
\l{profanacar. (\sou{trans., per)}. Violacar la santeso di la sakraji.}
\l{\cel{3}-(metaf.) Violacar la respekto debata ad ulu od ulo. DEFIRS}
\l{}
\l{profazo. (\sou{biol.}) Fazo unesma di mitoso.}
\l{}
\l{\sou{profermento.}\cel{2}(\sou{biol.)}\cel{3}Granularo fermentala qua, por divenar}
\l{\cel{3}efikiva e plear sua rolo zimazala, bezonas la prezenteso}
\l{\cel{3}di substanco preparanta.}
\l{}
\l{\sou{profesar}.\cel{2}(\sou{trans.)}\cel{2}Deklarar publike (sua fido religiala,}
\l{\cel{3}sua ateisteso, sua aprobo di...). - EFIS.}
\l{}
\l{profesiono. Sorto di agado per qua onu ganas sua vivo-neces-}
\l{\cel{3}aji. - DEFIRS.}
\l{}
\l{profesoro. Persono qua docas arto, cienco, en skolo superiora}
\l{\cel{3}(liceo, fakultato di yurocienco, di medicino). - DEFIRS.}
\l{}
\l{profeto. I. (\sou{religio}).La persono quan Deo inspiras,qua pre-}
\l{\cel{3}savigas la eventontaji, o revelas fakti celita de la per-}
\l{\cel{3}soni cetera. - II. La persono qua presavigas la eventont-}
\l{\cel{3}aji. - DEFIS.}
\l{}
\l{profilo. I. Aspekto di vizajo qua videsas de latere. - II.}
\l{\cel{3}Profilo di edifico, di fortreso, di tereno, di muluro:}
\l{\cel{3}(a) lia aspekto kande vidata de latere ; (b) lia seciono}
\l{\cel{3}sur plano, segun perpendiklo. - DEFIRS.}
}
